Many important contributions have been made to address the practical challenges of applying Shapley-based explanations. One such challenge is that standard SHAP implementations assume the input features are statistically independent, a condition rarely met in practice. When features are not independent, the model is evaluated on highly unrealistic input combinations. Another such challenge is that SHAP does not extend naturally to time-series data, as standard SHAP treats each feature–timestep pair as an independent player, causing the coalition space to grow exponentially and breaking temporal coherence. We review works that remedy these issues in turn.
\subsection{Manifold-Constrained SHAP}
Several works have addressed the off-manifold SHAP problem. Rather than assuming independence between features, \cite{aas2021explaining} propose three methods for modeling feature dependence: a Gaussian distribution, a Gaussian copula distribution, and an empirical conditional approach. This work is extended to non-parametric vine copulas \citep{aas2021vineexplaining}, and to mixed data with continuous and categorical features \citep{redelmeier2020explaining} by modeling feature dependence using conditional inference trees. To combat the computational complexity of SHAP and the additional complexity of estimating dependencies between all subsets of features, SHAP was adapted to groups of related features \citep{jullum2021groupshapley}. \citet{frye2020shapley} develop two solutions to on-manifold Shapley explainability: one based on generative modelling, which provides flexible access to data imputations, and another that directly learns the Shapley value function via supervised learning. \citet{olsen2022using} propose using a variational autoencoder with arbitrary conditioning \citep{ivanov2018variational}, building on the variational autoencoder framework \citep{kingma2014auto}, to model all feature dependencies simultaneously with a single model, reducing the computational burden of fitting separate models for each feature combination. A recent alternative approach produces Aumann-Shapley-type attributions via optimal generative flows, grounding baseline selection in optimal transport theory \citep{zhang2026axiomatic}. While effective at addressing feature dependence, none of these methods extend to temporal data.
\subsection{Temporal SHAP}
Concurrently but independently, works have extended SHAP to time-series data, though constrained to off-manifold implementations. TimeSHAP \citep{timeshap} extends SHAP to recurrent models by computing feature-, timestep-, and cell-level attributions through sequence perturbations. To manage the computational cost of long sequences, it introduces a pruning mechanism that discards early time steps whose cumulative relevance falls below a set threshold, under the assumption that the most recent time steps are most important to the prediction. WindowSHAP \citep{nayebi2023windowshap} combines neighboring time steps into windows and computes Shapley values at the window level, using stationary, sliding, and dynamic windowing strategies. \citet{villani2022feature} replace KernelSHAP's linear surrogate with a vector autoregressive model and introduce time-consistent Shapley values that respect temporal structure. ShaTS \citep{de2025shats} introduces a priori feature grouping strategies—temporal, sensor-level, and process-level—that preserve temporal dependencies during Shapley value computation rather than relying on post-hoc aggregation of individual feature–timestep attributions. GroupSegment-SHAP \citep{kim2026groupsegment} constructs group-segment players by partitioning the variable space via nonlinear dependence measures and detecting temporal distribution shifts, enabling attribution at the level of jointly active feature groups over contiguous time intervals.  All of these approaches are off-manifold.
\subsection{SHAP for Human Motion}
Literature applying SHAP to human motion is sparse. Preliminary works evaluate standard SHAP on motion datasets \citep{tempel2025choose, auletta2023predicting}. \citet{tempel2025explaining} design ShapGCN, which computes feature importance for a GCN model with human skeleton data structured as a graph, and validates explanations by perturbing the model's edge importance matrix to evaluate the impact of specific body keypoints on prediction outcomes. This perturbation strategy is specific to GCN architectures and therefore not applicable to other model types. Hence, temporal on-manifold SHAP for human motion has yet to be defined.


%The quality of these explanations is often evaluated using the principle of $\emph{faithfulness}$ \citep{lyu2024faithfulmodelexplanationnlp, dasgupta2022frameworkevaluatingfaithfulnesslocal, jacovi-goldberg-2020-towards}, which measures how accurately an explanation represents the reasoning of the underlying model. Two key aspects of faithfulness are $\emph{sufficiency}$ and $\emph{comprehensiveness}$ \citep{deyoung-etal-2020-eraser, yin2022sensitivitystabilitymodelinterpretations}; the former assesses whether the inputs deemed important are adequate for the model's prediction, and the latter examines if these features capture the essence of the model's decision-making process. We selected these metrics as they are widely-adopted, model-agnostic measures that directly assess explanation faithfulness through standard perturbation-based evaluation \citep{serrano-smith-2019-attention}, aligning with established principles of correctness and completeness in the explainability literature \citep{Nauta_2023}. 
